\section*{Abstract}
		Diese Arbeit stellt eine verhältnisbasierte Formulierung der T0-Physik vor, die mit deutlich weniger experimentellen Parametern auskommen soll. Aufbauend auf der vereinfachten Dirac-Gleichung und der universellen Lagrange-Funktion argumentieren wir, dass sich fundamentale Physik durch dimensionslose Energie-Skalen-Verhältnisse beschreiben lässt statt durch zugewiesene Parameter. Das T0-System benötigt einen SI-Referenzwert, um die verhältnisbasierte Physik mit messbaren Größen zu verbinden. In natürlichen Einheiten wird Einsteins $E = mc^2$ zur Identität $E = m$; die daraus folgenden Energie-Beziehungen gelten per Definition exakt, während zusammengesetzte Multi-Parameter-Formeln numerisch auf 99.98\% übereinstimmen. Die Physik soll sich auf Energie-Skalen-Verhältnisse zurückführen lassen, beschrieben durch die Feldgleichung $\partial^2 \Efield = 0$, mit quantitativen Vorhersagen über den einzigen Parameter $\xipar$ als SI-Referenzmaßstab (plus ganzzahlige Strukturwahlen als motivierte Setzungen).

	\medskip\noindent\textbf{Statusmarker.} Aussagen mit \textbf{[S]} sind \emph{spekulativ}: ein Hinweis oder eine Vermutung, die im Korpus noch nicht hergeleitet oder numerisch geprüft ist.

	\medskip\noindent\textbf{Hinweis zum Stand Myon $g-2$ (2025).} Dieses Dokument bezieht sich auf die Abweichung zwischen gemessenem $a_\mu$ und der Standardmodell-Vorhersage, wie sie 2021 vorlag. Mit dem Fermilab-Endergebnis und dem White Paper~2025 der Muon~$g-2$ Theory Initiative ist diese Abweichung nicht mehr signifikant. Aussagen, die die Anomalie als Beleg oder als Zielgröße eines T0-Beitrags verwenden, sind überholt; den aktuellen Stand und die Einordnung gibt Dok.~382.


	
	
	\section{Der T0-Ansatz: Von Parametern zu Verhältnissen}
	
	\subsection{Der Perspektivwechsel}
	
	Der T0-Ansatz schlägt einen Perspektivwechsel im Verständnis der Grundlagenphysik vor:
	
	\begin{tcolorbox}[colback=red!5!white,colframe=red!75!black,title=Gegenüberstellung]
		\textbf{Traditionelle Physik}: Multiple experimentelle Parameter
		\begin{itemize}
			\item $G = 6.67 \times 10^{-11}$ m³/(kg·s²) (gemessen)
			\item $\alpha = 1/137$ (gemessen)
			\item $m_e = 9.109 \times 10^{-31}$ kg (gemessen)
			\item 20+ unabhängige Parameter erforderlich
		\end{itemize}
		
		\textbf{T0-Verhältnis-basierte Physik}: Dimensionslose Skalen-Beziehungen
		\begin{itemize}
			\item Physik durch Energie-Skalen-Verhältnisse
			\item Ein SI-Referenzwert für quantitative Vorhersagen
			\item Mathematische Beziehungen, nicht experimentelle Parameter
			\item Reine Energie-Identitäten: $E = m$, $E = 1/L$, $E = 1/T$
		\end{itemize}
	\end{tcolorbox}
	
	\subsection{Aufbau auf T0-Grundlagen}
	
	Diese Arbeit bildet die dritte Stufe einer dreistufigen Entwicklung:
	
	\textbf{Stufe 1 - Vereinfachter Dirac}: Komplexe 4×4-Matrizen → Einfache Felddynamik $\partial^2 \deltam = 0$
	
	\textbf{Stufe 2 - Universelle Lagrange-Funktion}: 20+ Felder → Eine Gleichung $\Lag = \varepsilon \cdot (\partial \deltam)^2$
	
	\textbf{Stufe 3 - Verhältnis-basierte Physik}: Multiple Parameter → Energie-Skalen-Verhältnisse + SI-Referenz
	
	\subsection{Die Energie-Identität}
	
	In natürlichen Einheiten ($\hbar = c = 1$) nimmt Einsteins Gleichung die einfache Form an:
	
	\begin{equation}
		\boxed{E = m}
		\label{eq:energy_mass_identity}
	\end{equation}
	
	In diesen Einheiten handelt es sich um eine \textbf{Identität}, nicht um eine Umwandlung: Masse und Energie werden als dieselbe physikalische Größe behandelt.
	
	\begin{tcolorbox}[colback=blue!5!white,colframe=blue!75!black,title=Universelle Energie-Beziehungen]
		\textbf{Energie-Identitätssystem}:
		\begin{align}
			E &= m \quad \text{(Masse ist Energie)} \\
			E &= T_{\text{temp}} \quad \text{(Temperatur ist Energie)} \\
			E &= \omega \quad \text{(Frequenz ist Energie)} \\
			E &= \frac{1}{L} \quad \text{(Länge ist inverse Energie)} \\
			E &= \frac{1}{T} \quad \text{(Zeit ist inverse Energie)}
		\end{align}
		
		\textbf{Genauigkeit}: exakt per Definition (Identitäten in natürlichen Einheiten)
		
		\textbf{Komplexe Formeln}: 99.98-100.04\% (Rundungsfehler akkumulieren)
		
		\textbf{Folgerung}: Weniger zusammengesetzte Formeln führen weniger Rundungsfehler ein.
	\end{tcolorbox}
	
	\section{Teil I: Reine Verhältnis-basierte Physik (ohne SI-Referenz)}
	
	\subsection{Universelle Energiefeld-Dynamik}
	
	Im T0-Rahmen sind alle Teilchen Energie-Anregungsmuster im universellen Feld $\Efield(x,t)$:
	
	\begin{equation}
		\boxed{\partial^2 \Efield = 0}
		\label{eq:universal_field_equation}
	\end{equation}
	
	\textbf{Anspruch}: Diese Klein-Gordon-Gleichung für Energie soll alle Teilchen beschreiben.
	
	\subsection{Universelle Energie-Lagrange-Funktion}
	
	\begin{equation}
		\boxed{\Lag = \varepsilon \cdot (\partial \Efield)^2}
		\label{eq:universal_lagrangian}
	\end{equation}
	
	wo $\varepsilon$ die Energie-Skalen-Kopplung repräsentiert (dimensionsloses Verhältnis).
	
	\subsection{Antienergie: Symmetrie}
	
	\begin{equation}
		\boxed{\Efield_{\text{Antiteilchen}} = -\Efield_{\text{Teilchen}}}
		\label{eq:energy_antisymmetry}
	\end{equation}
	
	\textbf{Physikalisches Bild}: Positive und negative Energie-Anregungen desselben Feldes.
	
	\textbf{Lagrange-Universalität}:
	\begin{align}
		\Lag[+\Efield] &= \varepsilon \cdot (\partial \Efield)^2 \\
		\Lag[-\Efield] &= \varepsilon \cdot (\partial \Efield)^2
	\end{align}
	
	Dieselbe Physik für Teilchen und Antiteilchen durch Quadrierung.
	
	\subsection{Reine Verhältnis-Vorhersagen (unabhängig vom Wert von $\xipar$)}
	
	\subsubsection{Universelle Lepton-Verhältnisse}
	
	\begin{equation}
		\boxed{\frac{a_e^{(T0)}}{a_{\mu}^{(T0)}} = 1}
		\label{eq:universal_lepton_ratio}
	\end{equation}
	
	\textbf{Physikalische Bedeutung}: Alle Leptonen erhalten identische Energie-Korrekturen.
	
	\subsubsection{Energie-Unabhängigkeits-Verhältnisse}
	
	\begin{equation}
		\boxed{\frac{\Delta\Gamma^{\mu}(E_1)}{\Delta\Gamma^{\mu}(E_2)} = 1}
		\label{eq:energy_independence_ratio}
	\end{equation}
	
	\textbf{Unterscheidendes Merkmal}: Im Gegensatz zu Standardmodell-laufenden Kopplungen.

	\section{Teil II: Quantitative Vorhersagen (SI-Referenz erforderlich)}
	
	\subsection{Die SI-Referenz-Skala}
	
	Um quantitative Vorhersagen zu machen, benötigt die T0-Physik eine Verbindung zum SI-System:
	
	\begin{tcolorbox}[colback=green!5!white,colframe=green!75!black,title=SI-Referenz-Skala]
		\textbf{Definition}: $\xipar$ ist ein dimensionsloses Energie-Skalen-Verhältnis und der einzige Parameter des Rahmens; es ist geometrisch motiviert, nicht an Messdaten angepasst.
		
		\textbf{Higgs-Energie-Verhältnis}:
		\begin{equation}
			\xipar = \frac{\lambda_h^2 v^2}{16\pi^3 E_h^2}
		\end{equation}
		
		\textbf{Geometrisches Energie-Verhältnis}:
		\begin{equation}
			\xipar = \frac{2\ell_P}{\lambda_C}
		\end{equation}
		
		\textbf{SI-Referenzwert}: $\xipar = 1.33 \times 10^{-4}$
		
		\textbf{Rolle}: Verbindet dimensionslose Verhältnisse mit SI-messbaren Größen
	\end{tcolorbox}
	
	\subsection{Quantitative Lepton-Vorhersagen}
	
	Mit der SI-Referenz-Skala:
	
	\begin{equation}
		a_{\ell}^{(T0)} = \frac{1}{2\pi} \times \xipar^2 \times \frac{1}{12}
		\label{eq:quantitative_lepton_correction}
	\end{equation}
	
	\textbf{Numerische Berechnung}:
	\begin{align}
		a_{\ell}^{(T0)} &= \frac{1}{2\pi} \times (1.33 \times 10^{-4})^2 \times \frac{1}{12} \\
		&= \frac{1}{6.283} \times 1.77 \times 10^{-8} \times 0.0833 \\
		&= 2.47 \times 10^{-10}
	\end{align}
	
	\begin{tcolorbox}[colback=blue!5!white,colframe=blue!75!black,title=Universelle Lepton-Vorhersage]
		\textbf{Elektron g-2}: $a_e^{(T0)} = 2.47 \times 10^{-10}$
		
		\textbf{Myon g-2}: $a_{\mu}^{(T0)} = 2.47 \times 10^{-10}$ (identisch)
		
		\textbf{Tau g-2}: $a_{\tau}^{(T0)} = 2.47 \times 10^{-10}$ (gleicher Wert)
		
		\textbf{Aktuelle Myon-Anomalie}: $\Delta a_{\mu} \approx 25 \times 10^{-10}$
		
		\textbf{T0-Beitrag}: $\sim 10\%$ der beobachteten Anomalie
	\end{tcolorbox}
	
	\subsection{Quantitative QED-Vorhersagen}
	
	\begin{equation}
		\frac{\Delta\Gamma^{\mu}}{\Gamma^{\mu}} = \xipar^2 = 1.77 \times 10^{-8}
		\label{eq:quantitative_qed_correction}
	\end{equation}
	
	\textbf{Energie-Unabhängigkeits-Verifikation}:
	\begin{table}[htbp]
		\centering
		\adjustbox{max width=\textwidth}{%
\begin{tabular}{lcc}
			\toprule
			\textbf{Energie-Skala} & \textbf{T0-Korrektur} & \textbf{Standardmodell} \\
			\midrule
			1 MeV & $1.77 \times 10^{-8}$ & Laufende $\alpha(E)$ \\
			1 GeV & $1.77 \times 10^{-8}$ & Laufende $\alpha(E)$ \\
			100 GeV & $1.77 \times 10^{-8}$ & Laufende $\alpha(E)$ \\
			1 TeV & $1.77 \times 10^{-8}$ & Laufende $\alpha(E)$ \\
			\bottomrule
		\end{tabular}%
}
		\caption{Energie-unabhängige T0-Korrekturen vs. Standardmodell}
	\end{table}

	\section{Experimentelle Verifikationsstrategie}
	
	\subsection{Reine Verhältnis-Tests (Keine SI-Referenz benötigt)}
	
	\textbf{Test 1 - Universelle Lepton-Verhältnisse}:
	\begin{itemize}
		\item Messe $a_e^{(T0)}/a_{\mu}^{(T0)} = 1$
		\item Unabhängig von absoluten Werten
		\item Testet Universalitätsprinzip direkt
	\end{itemize}
	
	\textbf{Test 2 - Energie-Unabhängigkeit}:
	\begin{itemize}
		\item Messe QED-Korrekturen bei verschiedenen Energien
		\item Verhältnis sollte konstant sein: $\Delta\Gamma(E_1)/\Delta\Gamma(E_2) = 1$
		\item Unterscheidet von Standardmodell-laufenden Kopplungen
	\end{itemize}
	
	\textbf{Test 3 - Wellenlängen-Verhältnisse} \textbf{[S]}:
	\begin{itemize}
		\item Multi-Wellenlängen-Beobachtungen derselben Objekte
		\item Teste $z(\lambda_1)/z(\lambda_2) = \lambda_2/\lambda_1$
		\item Unabhängig von absoluter Rotverschiebungs-Kalibrierung
	\end{itemize}
	
	\subsection{Quantitative Tests (Erfordern SI-Referenz)}
	
	\textbf{Präzisions-g-2-Messungen}:
	\begin{itemize}
		\item Elektron g-2: Detektiere $2.47 \times 10^{-10}$ Korrektur
		\item Myon g-2: Prüfe $\sim 10\%$ der aktuellen Anomalie
		\item Tau g-2: Erste Messung, erwarte denselben Wert
	\end{itemize}
	
	\textbf{Multi-Energie-QED-Tests}:
	\begin{itemize}
		\item Messe absolut $\Delta\Gamma/\Gamma = 1.77 \times 10^{-8}$
		\item Verifiziere Energie-Unabhängigkeit über Dekaden
		\item Vergleiche mit Standardmodell-Vorhersagen
	\end{itemize}
	
	\section{Dunkle Materie und Dunkle Energie\\ aus Energie-Verhältnissen}
	
	\textbf{[S]} Die Aussagen dieses Abschnitts sind spekulativ: Der heutige Korpus klammert den kosmischen Sektor bewusst aus, weil auch das Standardmodell ihn nicht herleitet (P39).
	
	\subsection{Dunkle Materie: Unterschwellen-Energie-Oszillationen}
	
	\textbf{Verhältnis-basierte Beschreibung}:
	\begin{equation}
		\frac{\Efield_{\text{dunkel}}}{\Efield_{\text{Schwelle}}} = \xipar \sqrt{\frac{\rho_{\text{lokal}}}{\rho_{\text{kritisch}}}}
	\end{equation}
	
	\textbf{Physikalischer Mechanismus}: Zufallsphasen-Energie-Oszillationen unter der Teilchen-Detektionsschwelle.
	
	\subsection{Dunkle Energie: Großskalige Energie-Gradienten}
	
	\textbf{Verhältnis-basierte Energiedichte}:
	\begin{equation}
		\frac{\rho_{\Lambda}}{\rho_{\text{kritisch}}} = \frac{1}{2} \xipar^2 \left(\frac{E_{\text{Planck}}}{L_{\text{Hubble}} \cdot E_{\text{Planck}}}\right)^2
	\end{equation}
	
	\textbf{Quantitative Vorhersage}: $\rho_{\Lambda} \approx 6 \times 10^{-30}$ g/cm$^3$ (in der Größenordnung des beobachteten Werts)
	
	\section{Philosophische Einordnung: Physik ohne eigenständige Materie}
	
	\subsection{Reine Energie-Realität}
	
	\begin{tcolorbox}[colback=purple!5!white,colframe=purple!75!black,title=Entmaterialisierte Beschreibung]
		\textbf{Traditionelle Sicht}: Materie, Energie, Kräfte, Raumzeit als separate Entitäten
		
		\textbf{T0-Sicht}: Nur Energie-Muster und ihre Verhältnisse
		
		\textbf{Was wir Teilchen nennen}: Lokalisierte Energie-Konzentrationen
		
		\textbf{Was wir Kräfte nennen}: Energie-Gradienten-Wechselwirkungen
		
		\textbf{Was wir Raumzeit nennen}: Energie-Muster-Substrat
		
		\textbf{Was wir Bewusstsein nennen} \textbf{[S]}: Selbstreferentielle Energie-Muster
		
		\textbf{Kernaussage}: Reine Energie-Beziehungen regiert von $\partial^2 \Efield = 0$
	\end{tcolorbox}
	
	\subsection{Von Komplexität zu Einfachheit}
	
	\textbf{Physik-Evolution}:
	\begin{enumerate}
		\item \textbf{Antik}: Vier Elemente
		\item \textbf{Klassisch}: Teilchen in Raumzeit
		\item \textbf{Modern}: Felder und Kräfte
		\item \textbf{Standardmodell}: 20+ Parameter
		\item \textbf{T0-Ansatz}: Energie-Verhältnisse + eine SI-Referenz
	\end{enumerate}
	
	\textbf{Ziel ist weitgehende Vereinfachung}: möglichst wenige fundamentale Annahmen.
	
	\subsection{Bewusstsein und Energie-Muster}
	
	\textbf{Eine weiterführende Frage}: Wenn alles Energie-Muster sind, was ist mit dem Bewusstsein?
	
	\textbf{[S] T0-Deutung}: Bewusstsein wäre ein sich selbst beobachtendes Energie-Muster. Wir sind temporäre Organisationen des universellen Energiefelds, die die Fähigkeit zur Selbstreferenz und subjektiven Erfahrung entwickelt haben.
	
	\section{Das Verhältnis-Physik-Erbe}
	
	\subsection{Angestrebte Ergebnisse}
	
	Der verhältnisbasierte T0-Ansatz verfolgt folgende Ziele; der jeweilige Stand ist in den Detaildokumenten festgehalten:
	
	\begin{enumerate}
		\item \textbf{Parameter reduziert}: 20+ → $\xipar$ als SI-Referenz (plus ganzzahlige Setzungen)
		\item \textbf{Kräfte gemeinsam beschrieben}: Durch Energie-Gradienten-Wechselwirkungen
		\item \textbf{Teilchenvielfalt geordnet}: Als Energie-Muster beschrieben
		\item \textbf{Antiteilchen gedeutet}: Negative Energie-Anregungen
		\item \textbf{Gravitation einbezogen}: Über Energie-Raumzeit-Kopplung
		\item \textbf{Dunkle Phänomene} \textbf{[S]}: Als Energiefeld-Effekte gedeutet (spekulativ, P39)
		\item \textbf{Exakte Identitäten}: In natürlichen Einheiten per Definition
		\item \textbf{Verhältnis-basierte Physik formuliert}: Reine Skalen-Beziehungen
	\end{enumerate}
	
	\subsection{Die Zweistufige Teststrategie}
	
	\textbf{Stufe 1 - Reine Verhältnisse} (unabhängig vom Wert von $\xipar$):
	\begin{itemize}
		\item Universelle Lepton-Korrektur-Verhältnisse
		\item Energie-unabhängige QED-Verhältnisse
		\item Wellenlängenabhängige Rotverschiebungs-Verhältnisse \textbf{[S]}
		\item Gravitations-Modifikations-Verhältnisse
	\end{itemize}
	
	\textbf{Stufe 2 - Quantitative Vorhersagen} (SI-Referenz):
	\begin{itemize}
		\item Absolute g-2-Korrekturen
		\item Absolute QED-Vertex-Modifikationen
		\item Absolute kosmologische Parameter \textbf{[S]}
		\item Absolute dunkle Materie/Energie-Dichten \textbf{[S]}
	\end{itemize}
	
	\subsection{Stand des Rahmens}
	
	\begin{tcolorbox}[colback=yellow!5!white,colframe=orange!75!black,title=Grundstruktur im T0-Rahmen]
		\textbf{Im T0-Rahmen reduziert sich die Grundstruktur auf wenige Bausteine}.
		
		\textbf{Die fundamentale Gleichung}: $\partial^2 \Efield = 0$
		
		\textbf{Die universellen Verhältnisse}: Energie-Skalen-Beziehungen
		
		\textbf{Die SI-Verbindung}: Eine Referenz-Skala $\xipar$
		
		\textbf{Alles andere}: Verschiedene Lösungen und Muster
		
		\textbf{Tiefere Ebene}: Ob es eine tiefere Ebene gibt, bleibt offen; im T0-Rahmen ist keine weitere vorgesehen
		
		\textbf{Zukünftige Arbeit}: Vor allem Anwendungen und Messungen, die den Rahmen prüfen
	\end{tcolorbox}
	
	\section{Schlussfolgerung: Das Verhältnis-basierte Universum}
	
	\subsection{Zusammenfassung der Struktur}
	
	Im T0-Rahmen wird die Physik durch Energie-Skalen-Verhältnisse beschrieben:
	
	\textbf{Ebene 1}: Dimensionslose Energie-Verhältnisse (unabhängig von der SI-Referenz)
	
	\textbf{Ebene 2}: Eine SI-Referenz-Skala (quantitative Vorhersagen)
	
	\textbf{Ebene 3}: Reine Energie-Muster regiert von $\partial^2 \Efield = 0$
	
	Die beobachteten und gemessenen Phänomene sollen sich aus dieser einfachen verhältnisbasierten Struktur ergeben.
	
	\subsection{Rückblick}
	
	Der Weg führt von der Vielzahl der Parameter der traditionellen Physik zu einer einfacheren, verhältnisbasierten Energie-Dynamik.
	
	\textbf{Die Kernaussage}: Nach dieser Auffassung liegt der Kern der Natur nicht in komplizierter Mathematik oder exotischen Phänomenen, sondern in einfachen Skalen-Beziehungen.
	
	\textbf{Ein Feld}. \textbf{Eine Gleichung}. \textbf{Energie-Verhältnisse}. \textbf{Eine SI-Referenz}.
	
	Alles Weitere wären dann verschiedene Muster und Verhältnisse der Energie, bis hin zu dem Muster, das wir menschliches Bewusstsein nennen \textbf{[S]}.
	
	\begin{equation}
		\boxed{\text{Realität} = \text{Energie-Verhältnisse in } \Efield(x,t)}
	\end{equation}
	
	\textbf{Zusammengefasst: Im T0-Rahmen wird das Universum als Gefüge von Energie-Verhältnissen beschrieben; geprüft wird dies über die oben genannten Tests.}
	
\begin{thebibliography}{99}
	
	\bibitem{pascher_simplified_dirac_2025}
	Pascher, J. (2025). \textit{Vereinfachte Dirac-Gleichung in der T0-Theorie: Von komplexen 4×4-Matrizen zu einfacher Feld-Knoten-Dynamik}. \\
	\href{https://github.com/jpascher/T0-Time-Mass-Duality/blob/main/2/pdf/050_diracVereinfacht_En.pdf}{https://github.com/jpascher/T0-Time-Mass-Duality/blob/main/2/pdf\\/050\_diracVereinfacht\_En.pdf}
	
	\bibitem{pascher_lagrangian_comparison_2025}
	Pascher, J. (2025). \textit{Einfache Lagrange-Revolution: Von Standardmodell-Komplexität zu T0-Eleganz}. \\
	\href{https://github.com/jpascher/T0-Time-Mass-Duality/blob/main/2/pdf/049_LagrandianVergleich_En.pdf}{https://github.com/jpascher/T0-Time-Mass-Duality/blob/main/2/pdf\\/049\_LagrandianVergleich\_En.pdf}
	
	\bibitem{pascher_verification_table_2025}
	Pascher, J. (2025). \textit{T0-Modell-Verifikation: Skalen-Verhältnis-basierte Berechnungen vs. CODATA/Experimentelle Werte}. \\
	\href{https://github.com/jpascher/T0-Time-Mass-Duality/blob/main/2/pdf/054_Elimination_Of_Mass_Dirac_Tabelle_En.pdf}{https://github.com/jpascher/T0-Time-Mass-Duality/blob/main/2/pdf/\\054\_Elimination\_Of\_Mass\_Dirac\_Tabelle\_En.pdf}
	
	\bibitem{einstein_mass_energy_1905}
	Einstein, A. (1905). \textit{Ist die Trägheit eines Körpers von seinem Energieinhalt abhängig?} Ann. Phys. \textbf{17}, 639--641.
	
	\bibitem{dirac_original_1928}
	Dirac, P. A. M. (1928). \textit{The Quantum Theory of the Electron}. Proc. R. Soc. London A \textbf{117}, 610.
	
	\bibitem{muong2_experiment_2021}
	Myon g-2 Kollaboration (2021). \textit{Messung des positiven Myon-anomalen magnetischen Moments auf 0.46 ppm}. Phys. Rev. Lett. \textbf{126}, 141801.
	
	\bibitem{higgs_mechanism_1964}
	Higgs, P. W. (1964). \textit{Gebrochene Symmetrien und die Massen von Eichbosonen}. Phys. Rev. Lett. \textbf{13}, 508--509.
	
	\bibitem{planck_collaboration_2020}
	Planck Kollaboration (2020). \textit{Planck 2018 Ergebnisse. VI. Kosmologische Parameter}. Astron. Astrophys. \textbf{641}, A6.
	
	\bibitem{weinberg_qft_1995}
	Weinberg, S. (1995). \textit{Die Quantentheorie der Felder, Band 1: Grundlagen}. Cambridge University Press.
	
	\bibitem{particle_data_group_2022}
	Teilchendaten-Gruppe (2022). \textit{Übersicht der Teilchenphysik}. Prog. Theor. Exp. Phys. \textbf{2022}, 083C01.
	
\end{thebibliography}